
\vsize 18cm
\hsize 15.4cm
\magnification=1200
\item{}
\vskip 1cm
\font\r=cmr10 scaled\magstep2

\centerline{\r INVARIANTS AND RELATED LIAPUNOV}
\medskip
\centerline{\r FUNCTIONS FOR DIFFERENCE EQUATIONS}
\bigskip
\centerline{by}
\bigskip
\centerline{M.R.S.Kulenovi\'c}

\centerline{Department of Mathematics}

\centerline{University of Rhode Island}

\centerline{Kingston, RI 02881}

\vskip 1.5cm
\centerline{\bf Abstract}
\vskip 1.5cm
Consider the difference equation
$${\bf x_{n+1}}={\bf f( x_n)},$$
where ${\bf x_n}$ is in $R^k$ and ${\bf f}: D\to D$ is continuous where $D \subset R^k$.
Suppose that $I:R^k\to R$ is a continuous invariant that is, $I({\bf f(x)})=I({\bf x})$
for every ${\bf x}\in D$. We will show that if $I$ attains an isolated minimum
or maximum value
at the  equilibrium (fixed) point ${\bf p}$  of this system, then there exists a
Liapunov function, namely  $\pm(I({\bf x}) - I({\bf p}))$ and so the equilibrium
${\bf p}$ is stable. This result is then applied to some difference equations
appearing in different fields of applications.

\vskip 4.5cm

\noindent{\bf Key Words:} invariants, stability, Lyness' equation

\noindent{\bf AMS 1991 Subject Classification: Primary-39A11, Secondary-39Al0}


\vfill\eject
\vglue 2cm

{\r 1. Introduction}
\bigskip
Consider Lyness' equation
$$x_{n+1}={{a + x_n}\over x_{n-1}},$$
where $a>0$ is a parameter and $x_{-1}>0, x_0>0$. Recently, this equation
attracted a lot of attention, see [KL] and [GJKL], and its stability has been
proved showing that it
can be transformed to an area preserving map, see [KLTT]. Also KAM theory has
been applied
to provide another proof of stability, see [KLTT]. However, no Liapunov function
has been found so far. Here we will construct the Liapunov function for this
equation, using the well known invariant:
$$I(x_n,x_{n-1})=\left(1+{1\over x_n}\right)\left(1+{1\over x_{n-1}}\right)(a+x_n+x_{n-1}).$$
In fact, we will show that the Liapunov function $V(x,y)$ is given by
$$V(x,y)=I(x,y)-I(p,p)=I(x,y)-{(p+1)^3\over p},$$
where
$$p= {{1\pm \sqrt{1+4a}}\over 2}$$
is the equilibrium of this equation.

One of the possible third order generalizations of Lyness' equation is Todd's
equation which has the form:
$$x_{n+1}={{a + x_n + x_{n-1}}\over x_{n-2}},$$
where $a>0$ is a parameter and $x_{-1}, x_0, x_1>0$.
Here we will construct the Liapunov function for this equation using well
known invariant:
$$I(x_n,x_{n-1},x_{n-2})=\left(1+{1\over x_n}\right)\left(1+{1\over x_{n-1}}\right)\left(1+{1\over x_{n-2}}\right)(a+x_n+x_{n-1}+x_{n-2}).$$
In fact, we will show that the Liapunov function $V(x,y,z)$ is given by
$$V(x,y,z)=I(x,y,z)-I(p,p,p)=I(x,y,z)-{(p+1)^4/p^2},$$
where $p$ is the equilibrium of this equation.
\bigskip
\vfill\eject
{\r 2. Main Result. }
\bigskip
For the sake of completeness, we provide the definition of
stability that will be used in this paper.
\medskip
{\bf DEFINITION} ([KP], p.171) Let {\bf u} belong to ${\bf R}^n$ and $r>0$. The open
ball centered at {\bf u} with radius $r$ is the set
$$B({\bf u},r) = \{{\bf v} \in {\bf R}^n: |{\bf v} - {\bf u}|<r\}.$$
Let {\bf p} be a fixed point of {\bf f}. Then {\bf p} is stable provided that,
given any ball B({\bf p},$\epsilon$), there is a ball B({\bf p},$\delta$) so that
if {\bf u} is in B({\bf p},$\epsilon$), then ${\bf f}^k ({\bf p })$ is in B({\bf p},$\epsilon$)
for $k=1,2,...$. Here ${\bf f}^k ({\bf p })$ denotes $k$-th iterate of ${\bf p}$.
\medskip
Our main result is the following
\bigskip
{\bf THEOREM.} Consider the  difference equation
$${\bf x_{n+1}}={\bf f( x_n)}$$
where ${\bf x_n}$ is in $R^k$ and ${\bf f}: D\to D$ is continuous, where
$D \subset R^k$. Suppose that $I:R^k\to R$ is a continuous invariant that is
$I({\bf f(x)})=I({\bf x})$
for every ${\bf x}\in D$. If $I$ attains an isolated local minimum or maximum value
at the equilibrium (fixed) point ${\bf p}$  of this system, then there exists a
Liapunov function equal to  $\pm(I({\bf x}) - I({\bf p}))$ and so the equilibrium
${\bf p}$ is stable.
\medskip
{\bf Proof.} Assume that $I$ attains a local isolated minimum at ${\bf p}$. Then there
exists some neighborhood $D_p$ of $\bf p$ such that $I({\bf x})>I({\bf p})$ for
every ${\bf x}\in D_p$. In view of this we have

1. $V({\bf p}) = I({\bf p})  - I({\bf p}) = 0.$

2. $V({\bf x}) = I({\bf x}) - I({\bf p}) > 0$ , for {\bf x} in $D_p $.

3. $(V {\bf f})({\bf x}) = I({\bf f({\bf x})}) - I({\bf p}) = I({\bf x}) - I({\bf p}) = V({\bf x}).$

where in 3. we use the fact that $I$ is an invariant. Thus, all three conditions
in [KP, p.177] are satisfied and so $V$ is Liapunov function. Consequently,
${\bf p}$ is stable.

The proof in the case where $I$ attains its maximum follows immediately from
the result that was just proved by observing the simple fact that negative of
an invariant is also an invariant.
\bigskip
{\bf REMARK 1.} This theorem shows that whenever certain
conditions are satisfied, there exists a Liapunov function which
is actually a special invariant of considered equation. As is well
known, see [GM] and [M] the invariant is not uniquely determined
but rather if one invariant is given then any appropriate function
(continuous, analytic) of this invariant is invariant itself.
Furthermore, the difference equation can have several
"independent" invariants (eg. two invariants neither of which can
be derived from the other through homeomorphic changes of
variable), see [BB]. As we will mention later we have developed
{\it Mathematica} package that can almost automatically find the
rational invariant for a given equation (if there is one) and
using our result here construct the corresponding Liapunov
function, if such function exists.
\medskip
{\bf REMARK 2.} An important question that could be raised here is
what are the implications of this theorem on the linearized
stability theory for this difference equation. In particular, is
it true that all the eigenvalues of the linearized system
(assuming that {\bf f} is differentiable) $${\bf y_{n+1}}={\bf
f'(p) y_n}$$ necessarily lie on the unit circle ? The long
standing conjecture is that the difference equation $${\bf
x_{n+1}}={\bf f( x_n)}$$ has an invariant if and only if the
corresponding linearized equation has an invariant, which in turn
is equivalent to the fact that the matrix ${\bf f'(p)}$ does not
have all eigenvalues inside or outside unit circle.


\bigskip
{\r 3. Applications.}
\bigskip
Here, we will give several equations with known invariants and using the above
result we will obtain the Liapunov function, and prove the stability of their
equilibria. We note that the positive definiteness of the Hessian at the critical
point is verified by checking that all the principal minors are positive.
\bigskip
{\bf 1.} Lyness' equation (see [KL] and [KLTT]):

$$x_{n+1}={{a + x_n}\over x_{n-1}},$$
where $a>0$ is a parameter and $x_1>0, x_0>0$.
Invariant:
$$I(x_n,x_{n-1})=\left(1+{1\over x_n}\right)\left(1+{1\over x_{n-1}}\right)(a+x_n+x_{n-1}).$$
Equilibrium: $p= {{1\pm \sqrt{1+4a}}\over 2}$ satisfies $p^2-p-a=0$.

The necessary conditions for the extremum give:
$${\partial I \over \partial x}\equiv \left(1+{1\over y}\right)\left(1-{{y+a}\over x^2}\right)=0,$$
$${\partial I \over \partial y}\equiv \left(1+{1\over x}\right)\left(1-{{x+a}\over y^2}\right)=0,$$
which leads to $x=y$ and to the equation $x^2-x-a=0$. This shows that the
critical points are exactly the equilibrium points. Now, we will check the Hessian
at a positive critical point $p_+= {{1+ \sqrt{1+4a}}\over 2}$, which will be
denoted as $p$

$$H=\left(\matrix{A&B\cr
                  B&C\cr}\right),$$
where
$$A={\partial^2 I\over \partial x^2}(p,p)=2{{(p+1)(p+a)}\over p^4}>0, $$

$$B={\partial^2 I\over {\partial x \partial y}}(p,p)={{a - 2p^2}\over p^4}, $$
and
$$C= A.$$
Now,
$$det\; H=AC - B^2 = {(3a + 2(a+1)p)(a+2(a+1)p + 4p^2)/p^8}>0,$$
which implies that the invariant $I$ attains its minimum at $(p,p)$. Thus,
$$\min \{I(x,y): (x,y)\in D\}=I(p,p)={{(p+1)^2 (a+2p)}\over p^2},$$
and so by our theorem $V(x,y)=I(x,y) - {{(p+1)^2 (a+2p)}\over p^2}$, which shows  that $p$
is stable.
\bigskip
{\bf 2.} Todd's equation (see [KL]):
$$x_{n+1}={{a + x_n + x_{n-1}}\over x_{n-2}},$$
where $a>0$ is a parameter, and $x_1>0, \; x_0>0, \; x_1>0$.

Invariant: 
$$I(x_n,x_{n-1},x_{n-2})=\left(1+{1\over x_n}\right)\left(1+{1\over x_{n-1}}\right)\left(1+{1\over x_{n-2}}\right)(a+x_n+x_{n-1}+x_{n-2}).$$
Equilibrium:\quad $p=1\pm\sqrt{1+a}$ satisfies $p^2-2p-a=0$.

The necessary conditions for the extremum give:
$${\partial I \over \partial x}\equiv \left(1+{1\over y}\right)\left(1+{1\over z}\right)\left(1-{{a+y+z}\over x^2}\right)=0,$$
$${\partial I \over \partial y}\equiv \left(1+{1\over x}\right)\left(1+{1\over z}\right)\left(1-{{a+x+z}\over y^2}\right)=0,$$
and
$${\partial I \over \partial z}\equiv \left(1+{1\over x}\right)\left(1+{1\over y}\right)\left(1-{{a+x+y}\over z^2}\right)=0.$$

This leads to $x=y=z$, and to the equation $x^2-2x-a=0$, which shows that the
critical points are exactly the equilibrium points. Now, we will check the Hessian
at a positive critical point $(p,p,p)$ where $p= 1+ \sqrt{1+a}$.

$$H=\left(\matrix{A&B&B\cr
                  B&A&B\cr
                  B&B&A\cr}\right),$$
where
$$A={\partial^2 I\over \partial x^2}(p,p) ={\partial^2 I\over \partial y^2}(p,p)={\partial^2 I\over \partial z^2}(p,p) =2{{(p+1)^2 (a+2p)}\over p^5}>0, $$
and
$$B={\partial^2 I\over {\partial x \partial y}}(p,p)=
{\partial^2 I\over {\partial y \partial z}}(p,p)=
{\partial^2 I\over {\partial x \partial z}}(p,p)={{(p+1)(a+p-2p^2)}\over p^5}. $$
Now,
$$det\; H_1=A >0,\quad det$$
$$H_2=det{\left(\matrix{A&B\cr
                                  B&A\cr}\right)}={{(p+1)^2 (a+(2a+3)p+6p^2)(3a+(2a+5)p+2p^2)}\over p^10}>0 ,$$
and
$$det H=det{\left(\matrix{A&B&B\cr
                  B&A&B\cr
                  B&B&A\cr}\right)} = (A-B)^2(A+2B)$$
$$= 2 {{ (p+1)^3(2a+(a+3)p)(a+(2a+3)p+6p^2)}\over{p^{15}}}>0,$$
which shows, that the invariant $I$ attains its minimum at $(p,p,p)$. Thus,
$$\min \{I(x,y,z): (x,y,z)\in D\}=I(p,p,p)={{(p+1)^3(a+3p)}\over p^3},$$
and so by our theorem $V(x,y,z)=I(x,y,z)-{{(p+1)^3(a+3p)}\over p^3}$, which shows, that
$p$ is stable.

\bigskip
{\bf 3.} $k+2$-th order Lyness' equation (see [KL]):
$$x_{n+1}={{a + x_n + x_{n-1}+...+x_{n-k}}\over x_{n-k-1}},$$
where $a>0$ is a parameter, and $x_-k>0,...,x_-1>0$.

Invariant:
$$I(x_n,...,x_{n-k-1})=\left(1+{1\over x_n}\right)\left(1+{1\over x_{n-1}}\right)...\left(1+{1\over x_{n-k-1}}\right)(a+x_n+x_{n-1}+...+x_{n-k-1}).$$

Equilibrium:\quad $p={{k+1\pm\sqrt{(k+1)^2+4a}}\over 2}$ satisfies $p^2-(k+1)p-a=0$.
The necessary conditions for the extremum give:
$${\partial I \over \partial x_{n-i}}\equiv \prod_{j=0,j\neq i}^{k+1}\left(1+{1\over x_{n-j}}\right)
\left(1-{{a+\sum_{j=0,j\neq i}^{k+1} x_{n-j}}\over x_{n-i}^2}\right)=0,\quad i=1,...,k+1,$$
which gives
$$x_{n-i}^2-a- \sum_{j=0, j\neq i}^{k+1} x_{n-j}=0,\quad i=1,...,k+1.$$
Subtracting the consecutive equations we get $x_n=x_{n-1}=...=x_{n-k-1}=p$,
which shows that the critical point is $(p,p,...,p)$.

Now, we will check the Hessian at a positive critical point $(p_+,...,p_+) $,
which  for simplicity will be denoted by $(p,...,p)$.
Computing the second order derivatives of $I$ with respect to $x_{n-i}$ we get
$${\partial^2 I\over \partial x_{n-i}^2}={2\over x_{n-i}^3}\left(a+\sum_{j=0, j\neq i}^{k+1} x_{n-j}\right)
\prod_{j=0,j\neq i}^{k+1}\left(1+{1\over x_{n-j}}\right),$$
and
$${\partial^2 I\over {\partial x_{n-i}}{\partial x_{n-j}}}=\prod_{p=0,p\neq i,j}^{k+1}\left(1+{1\over x_{n-p}}\right)
\left({{a+\sum_{p=0, p\neq i,j}^{k+1} x_{n-p}}\over{x_{n-i}^2 x_{n-j}^2}}-{1\over x_{n-i}^2}-{1\over x_{n-j}^2} \right),$$
$i=1,...,k+1$ at the equilibrium $(p,p,...,p)$, we get
$$A={\partial^2 I\over \partial x_{n-i}^2}|(p,p,...,p)=2{{(p+1)^{k+1} (a+(k+1)p)}\over p^{k+4}}>0,$$
and
$$B={\partial^2 I\over {\partial x_{n-i}}{\partial x_{n-j}}}|(p,p,...,p)={{(p+1)^{k} (a+kp-2p^2)}\over p^{k+4}}.$$
Thus, the Hessian takes the form
$$H=\pmatrix{A&B&\ldots&B\cr
             B&A&\ldots&B\cr
             \vdots&\vdots&\ddots&\vdots\cr
             B&B&\ldots&A\cr}.$$

Now,
$$det\; H_1=A >0,$$
$$det\;H_m=det{\pmatrix{A&B&\ldots&B\cr
                                             B&A&\ldots&B\cr
                                             \vdots&\vdots&\ddots&\vdots\cr
                                             B&B&\ldots&A\cr}}=(A-B)^{m+1}[A+(m-1)B]$$

$$={(p+1)^{(m+2)k}\over p^{(k+4)(m+2)}} [2(k+2)p^2 + (k+2+a)p + a]^{m+1} [(2(k-m)+4)p^2 + (2a+(m+1)k +2)p +(m+1)a]>0,$$  $k=1,...,m,$ where $H_m$ is m-th order principal minor of the Hessian $H$.

This shows, that the invariant $I$ attains its minimum at $(p,p,...,p)$. Thus,
$$\min \{I(x_n,x_{n-1},...,x_{n-k-1}):(x_n,x_{n-1},...,x_{n-k-1}) \in D\}=I(p,p,...,p)$$
$$={{(p+1)^{k+2} (a+(k+2)p)}\over p^{k+2}}$$
and so, by our theorem   

$$V(x_n,x_{n-1},...,x_{n-k-1})=I(x_n,x_{n-1},...,x_{n-k-1})-{{(p+1)^{k+2} (a+(k+2)p)}\over p^{k+2}},$$
which shows, that $p$ is stable.

Obviously, for $k=0$ and $k=1$ we get Lyness' and Todd's equation
respectively, and so this example generalizes the examples 1. and
2. Another generalization of Lyness' equation is given in the
following example.
\bigskip
{\bf 4.} Generalized Lyness' equation (see [FJL], [GJKL] and [KL]):
$$x_{n+1}={{a x_n + b}\over {(cx_n+d)x_{n-1}}},$$
where $a,b,c,d$ are positive parameters and $x_{-1}>0, x_0>0$.

Invariant:
$$I(x_n,x_{n-1})={{ab}\over {x_{n-1}x_n}} + (a^2+bd)\left({1\over x_n} +{1\over x_{n-1}}\right)$$
$$ + ad\left({x_{n-1}\over x_n} + {x_{n}\over x_{n-1}}\right) +
(ac + d^2)(x_n+x_{n-1}) + cdx_{n-1}x_n. $$
Equilibrium: $p$ satisfies  $cp^3+dp^2-ap-b=0$. Obviously, this equation has at
least one positive solution and we will investigate the stability of this
equilibrium.
The necessary conditions for the extremum give:
$${\partial I \over \partial x}\equiv -{{ab}\over {x^2y}} + cdy - {{a^2+bd}\over x^2} + ad\left({1\over y} - {y\over x^2}\right) +ac + d^2=0,$$
and
$${\partial I \over \partial y}\equiv -{{ab}\over {xy^2}} + cdx - {{a^2+bd}\over y^2} + ad\left({1\over x} - {x\over y^2}\right) +ac + d^2 =0,$$
which after some manipulations leads to $x=y$ and to the equation
$$cdx^4+(ac+d^2)x^3-(a^2+bd)x-ab= (dx+a)(cx^3+dx^2-ax-b)=0.$$
This shows that the critical points are exactly the equilibrium points. Now, we
will check the Hessian at a positive critical point $(p_+,p_+)$ which will be
denoted simply as $(p,p)$:
$$H=\left(\matrix{A&B\cr
                  B&A\cr}\right),$$
where
$$A={\partial^2 I\over \partial x^2}(p,p)={2\over p^4}[adp^2+(a^2+bd)p+ab]>0,$$
and
$$B={\partial^2 I\over {\partial x \partial y}}(p,p)={{cdp^4-2adp^2+ab}\over p^4}. $$
Now,
$$det\; H=A^2 - B^2 = (A-B)(A+B)={1\over p^8}[cdp^4+2(a^2+bd)p+3ab]E,$$
where by using the equilibrium equation
$$E=-cdp^4+4adp^2+2(a^2+bd)p+ab=d^2p^3+3adp^2+2a^2p+ab>0,$$
which shows, that the Hessian $H$ is positive definite at the point $(p,p),$
which means that the invariant $I$ attains its minimum at $(p,p)$. Thus,
$$\min \{I(x,y): (x,y)\in D\}=I(p,p),$$
and so by our theorem $V(x,y)=I(x,y)-I(p,p)$ which shows, that $p$ is stable.
\bigskip
{\bf 5.}  Gumovski-Mira equation (see [GM],[M] and [L]):
$$x_{n+1}={{2ax_n}\over {1+x_n^2}} - x_{n-1},$$
where $a>1$ is a parameter.

Invariant: $I(x_n,x_{n-1})={x_n}^2 {x_{n-1}}^2 + {x_n}^2 + {x_{n-1}}^2 - 2ax_n x_{n-1}$.

Equilibrium: $p=\sqrt{a-1}>0$ satisfies  $p^2=a-1$.

The necessary conditions for the extremum give:
$${\partial I \over \partial x}\equiv 2xy^2+2x-2ay=0,$$
and
$${\partial I \over \partial y}\equiv 2x^2y+2y-2ax=0,$$
which immediately leads to $x=y=p$.
This shows that the critical point is exactly the equilibrium point.
Now, we will check the Hessian at a positive critical point $(p,p)$.
$$H=\left(\matrix{A&B\cr
                  B&A\cr}\right),$$
where
$$A={\partial^2 I\over \partial x^2}(p,p)=2p^2+2>0,$$
and
$$B={\partial^2 I\over {\partial x \partial y}}(p,p)=4p^2-2a. $$
Now, by using the equilibrium equation we get
$$det\; H=A^2 - B^2 = (A-B)(A+B)=4(a+1-p^2)(3p^2+1-a)=16(a-1)>0.$$
This shows, that the Hessian $H$ is positive definite at the point $(p,p)$
which means that the invariant $I$ attains its minimum at $(p,p)$. Thus,
$$\min \{I(x,y): (x,y)\in D\}=I(p,p)=-(a-1)^2,$$
and so by our theorem $V(x,y)=I(x,y)+(a-1)^2,$ which shows, that $p$ is stable.
\bigskip
{\bf REMARK 3.} It is worth mentioning that all equations with
their invariants were known, and that we have proved the stability
of the equilibrium point using our result. The result that we have
proved emphasizes the importance of the invariants and of finding
some general methods for their construction. Especially we are
interested in automatic construction of the invariants and
corresponding Liapunov functions, (if they exist) that can be
implemented by using some of CAS (Computer Algebra System). In
this regard, we have developed {\it Mathematica} package, see [MK]
which was successfully tested on all applications of this paper.
More applications of this package will be presented in the coming
communications.

{\bf ACKNOWLEDGMENT.} The author is grateful to the referee for some helpful
comments.

\vfill\eject

\centerline{\bf REFERENCES}
\vskip 2cm
\item{[BB]} {{\bf J.M. Borwein and P.B. Borwein,} {\it Pi and the AGM,}
Wiley, 1987.}
\bigskip
\item{[FJL]} {{\bf J.Feuer, E.Janowski and G.Ladas,} Invariants for Some
Rational Recursive Sequences with Periodic Coefficients, {\it  J. Diff. Equa. Appl.} 2(1996), 167-174.}
\bigskip
\item{[GJKL]} {{\bf E.A.Grove, E.J.Janowski, C.M.Kent and G. Ladas,} On
the rational recursive sequence  $x_{n+1}=(\alpha x_n + \beta)/((\gamma x_n+\delta)x_{n-1})$,
{\it Comm. Applied Nonlinear Anal.} 1(1994), 61-72.}
\bigskip
\item{[GM]} {{\bf I.Gumovski and C.Mira,} {\it Recurrences and discrete dynamic
systems,} Lecture Notes in Mathematics, Springer, 1980.}
\bigskip
\item{[KL]}  {{\bf V.L.Kocic and G.Ladas,} {\it Global Behavior of Nonlinear
Difference Equations of Higher Order with Applications,} Kluwer Academic
Publishers, 1993.}
\bigskip
\item{[KLTT]}  {{\bf V.L.Kocic, G.Ladas, G.Tzanetopoulos and E.Thomas,} On
the Stability of Lyness' Equation, {\it Dynam. Contin. Discrete Impuls. Systems}, 1(1995),245-254.}
\bigskip
\item{[KP]} {{\bf W.G.Kelley and A.C.Peterson,} {\it Difference Equations,}
Academic Press, 1991}
\bigskip
\item{[L]}  {{\bf H.Lauwerier,} {\it Fractals,} Princeton University Press, 1991.}
\bigskip
\item{[M]}  {{\bf C.Mira,} {\it Chaotic Dynamics,} World Scientific, 1987.}
\bigskip
\item{[MK]} {{\bf O.Merino and M.R.S.Kulenovi\'c,} IFL-Invariants and Liapunov
Functions for Difference Equations in Mathematica, (to appear).}




\end
